
Three experiments are designed to answer three research questions mapping to the three-step Goodhart saturation mechanism:

\textbf{RQ1:} Is there a statistically significant structural break in PwC benchmark residual CoV at some paper\_count threshold? \textit{(H-E1)}

\textbf{RQ2:} Does the pre-breakpoint benchmark segment exhibit substantially higher performance variance than the global baseline? \textit{(H-M1)}

\textbf{RQ3:} Does the post-breakpoint segment exhibit substantially lower variance than the pre-breakpoint segment? \textit{(H-M2, H-M3)}

\subsubsection{4.1 Dataset}

PwC benchmark archive, August 2026 snapshot, N=115 benchmarks after min\_papers=38 filtering, paper\_count $\in$ [38, 352].

\subsubsection{4.2 Baselines}

\begin{table}[htbp]
\centering
\resizebox{\columnwidth}{!}{%
\begin{tabular}{lll}
\toprule
Method & Description & Why Included \\
\midrule
OLS trend & Linear regression CoV ~ paper\_count (rho=0.137) & Tests null: saturation is a smooth trend with no structural break \\
S\_index [Akhtar2026] & Composite LLM saturation metric & Closest existing saturation index --- contrasted against the threshold approach \\
Liao et al. aggregate CoV [Liao2022] & Time-based saturation metric & Largest saturation study --- paper\_count threshold contrasted against time-based aggregate \\
\bottomrule
\end{tabular}
}
\end{table}

\subsubsection{4.3 Evaluation Metrics}

\textbf{RQ1:} Permutation p (primary gate, < 0.05), paper\_count\* location (gate $\in$ [10, 120]), piecewise F-test p (corroborative).

\textbf{RQ2:} Pre-segment variance ratio > 1.0; F-test one-tailed p < 0.10.

\textbf{RQ3:} Brown-Forsythe variance ratio < 1.0, p < 0.05 (H-M2); $\geq$ 2/4 directional metrics at p < 0.10 (H-M3).

